\section{Related Work}
\label{sec:related}

\subsection{Switching Strategies}

Three main switching strategies are found in the NoC literature.
\emph{Store-and-forward} switching requires a complete packet to be buffered at
each router before forwarding, providing clean error isolation but imposing large
buffer requirements and high latency~\cite{dally2001}.
\emph{Wormhole switching}~\cite{dally1987} pipelines flits across routers,
dramatically reducing buffer depth requirements; latency depends on path length
and contention.
\emph{Circuit switching}~\cite{kumar2002} reserves a dedicated path before
transmission, eliminating mid-network contention at the cost of reservation
overhead.

Hybrid approaches have been explored. Æthereal~\cite{goossens2005} combines
guaranteed-throughput circuit-switched slots with best-effort wormhole traffic on
the same physical links. HyNoC differs: it uses source routing to establish a
dedicated path through a series of request/grant handshakes (the circuit phase),
then streams flits in wormhole fashion along that path (the data phase), with no
time-division multiplexing.

\subsection{Routing Algorithms}

XY deterministic routing~\cite{dally2001} is simple and deadlock-free on 2D
meshes but provides only a single path between any pair of nodes. Adaptive
routing~\cite{duato1993} improves load balancing but requires more complex router
logic and typically needs virtual channels to avoid deadlock.

\emph{Source routing} encodes the complete route in the packet header at the
sender. This approach, studied for NoCs in~\cite{liwei2007} and~\cite{mubeen2010},
shifts routing complexity from the router to the sending node or to an offline
compilation step. The router itself becomes simpler: it only needs to decode the
next hop from the current flit, forward it, and update the index. Source routing
naturally supports any topology without per-router routing table storage, and
allows the sender to choose among multiple paths.

HyNoC adopts source routing with a compact \emph{relative} hop encoding: since a
packet cannot exit from the same port it entered, a router with $P$ ports offers
only $P-1$ egress choices, encoded in $\lceil\log_2(P-1)\rceil$ bits per hop.
This reduces both header overhead and crossbar complexity simultaneously.

\subsection{Virtual Channels}

Virtual channels (VCs), introduced by Dally~\cite{dally1992}, multiplex several
logical channels over a single physical link, providing additional routing paths
and breaking deadlock cycles. Their main drawback is area: each VC requires a
dedicated buffer, and~\cite{mello2005} shows that router area scales nearly
linearly with VC count. On FPGAs, where block RAMs are shared between the
application and the communication infrastructure, this overhead is particularly
significant.

Several works have explored VC reduction or elimination.
Turn-model routing~\cite{glass1992} achieves deadlock freedom without VCs by
restricting turns in the routing graph. HyNoC's approach---providing path diversity
through richer topologies and compile-time route selection---aligns with this
trend, and is analyzed quantitatively in Section~\ref{sec:switching}.

\subsection{NoC Arbitration}

Round-robin arbiters are widely used for their fairness guarantees. Sequential
round-robin FSMs scan requests cyclically, introducing latency proportional to the
number of ports. HyNoC implements a \emph{parallel round-robin arbiter} (PRRA)
that combines the fairness of round-robin with a low-latency response: the arbiter
transitions directly to the state of any pending request, respecting a rotating
priority order, and responds in one or two cycles regardless of the number of
ports. The implementation relies on a LUT-mux structure that keeps the critical
path independent of port count.

\subsection{FPGA-Targeted NoCs}

Several NoC designs target FPGAs specifically. CONNECT~\cite{papamichael2012}
provides a parameterized, synthesizable NoC generator. HopliteRT~\cite{nair2019}
proposes a lightweight deflection-routing NoC without input buffers, achieving
extremely small area at the cost of non-deterministic latency. HyNoC targets a
different point: a fully parameterized, topology-agnostic NoC with bounded latency,
suited to predictable VLIW workloads.

\subsection{VLIW Many-Core Systems}

VLIW architectures expose instruction-level parallelism explicitly, making them
well suited for signal processing and scientific computing kernels~\cite{ellis1986}.
Distributed VLIW many-core systems---exemplified by the Kalray
MPPA~\cite{dinechin2013}---replicate simple VLIW cores connected through a NoC. In
such systems, communication patterns can be determined at compile time, making
static source routing particularly attractive.
